\begin{abstract}
Near-critical black-hole light paths produce image-order channels with different spatial maps, transfer weights, and distributed delays. We study which components of an earlier source history these channels support, separating historical reach, historical dimension, and held-out recovery. For separable source spaces, temporal factors whose support misses the direct retarded-time footprint generate exactly zero direct columns; added orders can make these columns nonzero without necessarily making them independent. In a twelve-geometry sampled-operator audit at matched direct-reference $\SNR=100$, the first two indirect orders extend the anchor-connected detectable span by $40M$, $24M$, and one $4M$ grid step at inclinations $20^\circ$, $50^\circ$, and $75^\circ$. At the reconstruction geometry, the direct operator annihilates a 72-dimensional old-contrast target. The order-labelled stack retains two operational combinations after all other 152 source coefficients are profiled out; the available row-reweighting bounds permit one or two under their stated assumptions. Across a separate three-anchor source-enrichment study, the supported fraction decreases while the absolute supported count increases. Two sealed inverse benchmarks address different historical objects: baseline-inclusive emissivity-level span increases from $48M$ to $80M$ in the successful regimes, and median paired morphology-error reductions are $0.133$ for truncated SVD and $0.164$ for ridge. Severe mismatch, reliable multi-feature recovery, and stable morphology intervals remain negative. These are finite-model results under ideal order labelling, not demonstrated continuum-acquisition accuracy. The full transfer quadrature remains unqualified, and an index-sum compression control does not establish a physical order-resolution advantage. The Mahakal phenomenon denotes this separation between access to earlier source times and the dimension and fidelity of the history that can actually be recovered.
\end{abstract}

\section{Introduction}\label{sec:intro}
Light from the vicinity of a black hole can reach a distant observer along substantially different null paths. Direct rays and increasingly bent higher-order rays sample different source positions and emission times. The resulting image hierarchy has established demagnification, rotation, and delay structure \cite{gralla,johnson}. For a time-dependent source, however, an image order is not a single delayed snapshot. Its rays carry a distribution of delays, and different orders have overlapping retarded-time footprints. This distinction matters when the forward model is used to infer a source history rather than only to render an image.

The existence of old arriving photons does not settle the inverse problem. A source perturbation may lie inside a sampled time window but project onto an almost invisible direction. Several old perturbations may give the same observation. A reconstruction may improve in average error without yielding a reliable interval of historical morphology. These possibilities motivate three questions:
\emph{Historical reach} asks which earlier source times are sampled and where a prescribed localized perturbation exceeds a sensitivity threshold. \emph{Historical dimension} counts distinguishable directions in a declared source space at finite noise, including when other source coefficients are unknown. \emph{Historical recovery} tests a fixed estimator on unseen histories under a declared loss.

These questions concern different observables. Neither a delay maximum nor an operator rank alone answers all three.

\begin{definition}[The Mahakal phenomenon]
We use the term for the separation between gravitational access to earlier source times and the finite-model dimension and fidelity of the history supported by the observations. Higher-order footprints can enter time regions missed by the direct footprint; resolving that history in a richer source representation can expose null or weak directions hidden by a coarser description. The relation is geometry-, noise-, and source-class-dependent. It is not a new propagation law, a statement that source enrichment always destroys injectivity, or a claim that the black hole contains a recoverable movie.
\end{definition}

The analysis proceeds from a conditional support result to finite-noise information and finally to held-out reconstruction. Its most direct new information result concerns a 72-dimensional non-axisymmetric old-source target inside a 224-dimensional localized model. At the reference geometry, the sampled direct operator is exactly zero on this target. The labelled stack retains two above-threshold combinations even when all 152 remaining source coefficients are treated as nuisance parameters. This is a stronger question than asking whether old-target columns become nonzero with the remainder fixed, but it is still much weaker than reconstructing a history.

The recovery experiments deliberately distinguish an emissivity-level target from a resolution-aware morphology target. The first yields a longer baseline-inclusive stable span but is dominated by the spatially constant level. The second yields material aggregate error reductions but neither reliable multi-feature recovery nor a nonzero stable morphology interval. Presenting these positive and negative endpoints together identifies what the likelihood and the estimators actually support.

\subsection{Scope and contribution}
The setting is a known, stationary Kerr or Schwarzschild geometry, an equatorial source, and a monochromatic scalar-intensity image-domain model. We retain orders $n=0,1,2$ under the ray-map classification. Multi-geometry sensitivity is evaluated on twelve prescribed spin--inclination pairs; source-class enrichment uses three prescribed anchors; held-out inversion uses one reference geometry. These are separate experimental populations.

The contribution is an organized information-and-recovery analysis for explicitly specified finite operators. It combines the support-nullity theorem, nuisance-adjusted old-target information, source-class stress, and two different held-out recovery criteria. Continuous transfer-integral accuracy and the interpretation of a physically unresolved image are not established by the archived computations. Throughout, numerical results are tied to their source representation, screen weights, noise law, and evaluation metric. The compiled experiment records and amendments are the source of the reported numbers \cite{archive}.

\section{Relation to existing work}\label{sec:related}
The highly bent Kerr image hierarchy and its long-baseline interferometric signatures are established results \cite{gralla,johnson}. Time-domain studies examine delayed image structure through glimmer and correlations \cite{wong,hadar}. These observables can encode both lensing properties and information about the emitting source. Our question is not whether delayed radiation or source-sensitive correlations exist, but which source directions are distinguishable under a specified measurement model and source representation.

AART provides the analytic Kerr ray-tracing framework underlying the archived maps \cite{aart}. Accuracy of the underlying geodesic expressions is distinct from accuracy of a downstream detector integral constructed by sampling, masking, interpolation, and quadrature. The present paper therefore separates pointwise transfer checks, discrete linear-algebra checks, and full acquisition accuracy.

Related inverse approaches address different targets and assumptions. Hotspot-based spacetime tomography studies recovery within a constrained source model \cite{tiede}. Dynamic interferometric imaging reconstructs time-varying images using an explicit temporal model \cite{bouman}. Orbital polarimetric tomography combines a gravitational model and an emission representation to infer flare structure \cite{levis}. Such work establishes important reconstruction precedents. Here the target geometry is fixed, the historical perturbation space is declared, and finite-model support and nuisance-adjusted information are evaluated separately from reconstruction.

The distinction between mathematical identifiability and stable inversion also belongs to the broader literature on light-ray transforms and discrete ill-posed problems \cite{lassas,hansen}. Regularization and prior assumptions can select or stabilize a solution, but the reported estimator must be distinguished from the information supplied by the likelihood \cite{stuart}. We make no literature-wide priority claim for that distinction. The specific object studied here is the retarded-time, image-order-labelled source operator and the bounded historical questions tested with it.
